\subsection{1840 Ceza Kanunnamesinin İçeriği}
\indent\indent 1840(1256) tarihli ceza kanunnamesi, mukaddime, on üç fasıl, kırk bir madde ve hatîmeden meydana gelmektedir. Mukaddimede açık bir şekilde Tanzimat Fermanı'na atıf bulunmkatadır. Mukaddime şöyle başlamaktır:
\begin{quote}
Beyaz üzerine şeref-sudür buyurulup geçen elli beş senesi Şa'ban-ı Şerifinin yirmi altıncı günü Gülhane'de kıraat olunan Hatt-ı Hümâyün-u adalet-makrün-i Hazret-i Şahâne mücebince kâffe-i tebea-i Devlet-i Aliyye bila istisna emniyet-i can ve mal ve mahfüziyet-i ırz ve namus hukuk-u mefrüzasına ez-ser nev-nail olmuş ve ber-muktezay-i hürriyet-i şer`iyye huzuru şer` ve kanunda ve mevadd-ı hukukiyede herkesin yeksan ve siyyan olması umür-i tabiiyyeden bulunmuş olmakla...\footfullcite[810-811]{islam_osmanli_hukuku} 
\end{quote}

Kanunnamedeki fasılları aşağıdaki şekilde özetleyebiliriz:\footcite[811-818]{islam_osmanli_hukuku}
\begin{itemize}
    \item Birinci Fasıl: Devlet aleyhine fitne ve katl suçları ve cezaları(kısas) üzerine dört madde
    \item İkinci Fasıl: Söz ve hareket ile karışıklık çıkarma suçu ve cezası(kürek) üzerine dört madde
    \item Üçüncü Fasıl: Irz ve namusa karşı işlenen suçlar(hakaret, darb vb.) ve cezaları(hapis) üzerine beş madde
    \item Dördüncü Fasıl: Mal ve emlâka karşı işlenen suçlar(hırsızlık, el koymak vb.) ve cezaları(sürgün) üzerine iki madde
    \item Beşinci Fasıl: Rüşvet suçu ve cezaları(tenzil-i rütbe, azil, kürek, sürgün) üzerine yedi madde
    \item Altıncı Fasıl: Devlet memurlarının işlediği hırsızlık suçu ve cezaları(tenzil-i rütbe, azil, kürek) üzerine üç madde
    \item Yedinci Fasıl: Devlet malının çalınması suçu ve cezaları(azil) üzerine dört madde
    \item Sekizinci Fasıl: Devlet görevlilerinin sorumluluklar ve anlaşmazlıklarında uygulanacak ceza(tedip) üzerine iki madde
    \item Dokuzuncu Fasıl: Verginin büyük küçük herkesin mâli kuvvetine göre toplanmasına, zaptiye ve raptiyeye muhalefet ve mukavemet suçu ve cezaları(hapis, kürek, kısas) üç madde
    \item Onuncu Fasıl: Silahlı yaralama ve öldürme suçu ve cezaları(kürek, kısas) üzerine iki madde
    \item On Birinci Fasıl: Yol kesme suçu ve cezaları(kürek, kısas) üzerine üç madde
    \item On İkinci Fasıl: Herkesin kanun önünde eşit olduğu ve kimsenin sebepsiz yere cezalandırılamayacağı üzerine iki madde
    \item On Üçüncü Fasıl: Taşra memurlarına muhalefet suçunun da işbu ceza kanunnamesinin kapsamına girmesi
\end{itemize}

Yukarıda da görüldüğü üzere, 1840 Ceza Kanunnamesi devlete ve toplum asayişine karşı işlenen suçları ve cezaları düzenlemiştir. Ceza türlerini genel olarak kısas, kürek, nefy(sürgün) ve hapis olarak dört ana kategoride toplamak mümkündür. 
Devlet memurlarınının işlediği suçlar tenzil-i rütbe, azil, tedip gibi daha hafif cezalar ile düzenlenmiştir. Hatîmede ise bir kere daha bu kanunnamenin \textit{Devlet-i Aliyyenin merâtib-i adide ve mütefavitesinde bulunanlar ve ehl-i islam ve milel-i saireden olanlar haklarında bila istisna vela garaz icra olunmak üzere karargir}
olduğu belirtilmiştir.\footcite[819]{islam_osmanli_hukuku}